\documentclass[11pt]{article}
\usepackage[utf8]{inputenc}
\usepackage[spanish]{babel}
\usepackage[margin=1in]{geometry}
\usepackage{amsmath,amssymb,booktabs,array,enumitem,xcolor,fancyhdr}
\usepackage{tikz}
\usetikzlibrary{babel}
\setlength{\parindent}{0pt}
\setlength{\parskip}{.45em}
\setlength{\headheight}{14pt}
\pagestyle{fancy}
\fancyhf{}\lhead{Descenso del gradiente}\rhead{Libreto de tablero}\cfoot{\thepage}
\title{\textbf{Descenso del gradiente con dos variables}\\[.3em]\large Libreto de tablero para una clase de 2 horas}
\author{Administración de Empresas --- enfoque intuitivo}
\date{}

\begin{document}
\maketitle

\section*{Propósito y agenda}
Este libreto indica qué escribir, qué preguntar y qué cuenta resolver. El cuaderno
contiene las dos animaciones incrustadas: una trayectoria en curvas de nivel y otra
sobre una superficie 3D.

\begin{center}
\begin{tabular}{p{1.5cm}p{5cm}p{6.3cm}}
\toprule
Min. & Bloque & Evidencia de aprendizaje\\
\midrule
0--10 & Problema de negocio & Reconoce dos decisiones simultáneas.\\
10--25 & Mapa de costos & Lee un punto como una política.\\
25--40 & Dos señales & Distingue $g_x$ y $g_y$.\\
40--68 & Caso y cuentas & Ejecuta tres rondas completas.\\
68--73 & Pausa & Explica la idea sin fórmulas.\\
73--90 & Animaciones & Interpreta trayectoria y altura.\\
90--105 & Tamaño del paso & Compara estabilidad.\\
105--115 & Comité empresarial & Resuelve un nuevo punto.\\
115--120 & Cierre & Formula la idea en sus palabras.\\
\bottomrule
\end{tabular}
\end{center}

\newpage
\section*{Tablero 1 --- Dos decisiones, un objetivo (0--10 min)}
\textbf{Diga:} ``Una cadena minorista decide cuánto invertir en publicidad y cuánto
inventario de seguridad mantener. Ambas decisiones afectan el costo. Hoy no elegiremos
una y congelaremos la otra: las corregiremos juntas.''

\textbf{Escriba en grande:}
\[
\boxed{\text{publicidad }x\qquad+\qquad\text{inventario }y
\qquad\longrightarrow\qquad\text{costo semanal}}
\]
\[
\boxed{\text{MEDIR}\longrightarrow\text{CORREGIR}\longrightarrow\text{REPETIR}}
\]

\textbf{Pregunte:} ``¿Qué puede pasar si hay poca publicidad? ¿Y si hay demasiado
inventario?'' No busque una fórmula todavía; recoja intuiciones.

\section*{Tablero 2 --- El paisaje de dos variables (10--25 min)}
\textbf{Dibuje una montaña vista desde arriba:} dos ejes, varios óvalos concéntricos
y un punto lejos del centro. Marque el centro como ``menor costo''.

\[
\begin{array}{c}
\text{un punto }(x,y)=\text{una política completa}\\
\text{curva de nivel}=\text{políticas con costo parecido}\\
\text{centro del mapa}=\text{valle o mejor zona}
\end{array}
\]

\textbf{Diga:} ``Con una variable caminábamos en una línea. Con dos variables
caminamos sobre un mapa. Cada paso puede cambiar publicidad e inventario al mismo
tiempo.''

\textbf{Pregunte:} ``Si estamos a la derecha del valle y arriba del valle, ¿qué
deberíamos corregir?'' Respuesta: ambas coordenadas, en las direcciones que bajan.

\newpage
\section*{Tablero 3 --- El caso y las dos señales (25--40 min)}
\textbf{Escriba el modelo y traduzca cada término:}
\[
\boxed{J(x,y)=(x-12)^2+2(y-8)^2+100}
\]
\[
\begin{array}{c|c}
x&\text{inversión semanal en publicidad}\\
y&\text{inventario de seguridad}\\
100&\text{costo base inevitable}\\
(x-12)^2,\ 2(y-8)^2&\text{costo por alejarse de la mejor zona}
\end{array}
\]

El equipo analítico entrega las señales:
\[
\boxed{g_x=2(x-12)}\qquad\boxed{g_y=4(y-8)}
\]

\textbf{Diga:} ``No vamos a demostrar estas expresiones. Las usaremos como un
reporte: $g_x$ indica cómo corregir publicidad y $g_y$ cómo corregir inventario.''

\textbf{Escriba la regla completa:}
\[
\boxed{x_{\rm nuevo}=x-\eta g_x,\qquad
y_{\rm nuevo}=y-\eta g_y}
\]
\[
\eta=\text{tamaño del paso: cuánto confiamos en cada corrección}.
\]

\section*{Derivadas parciales sin miedo (30--40 min)}
\textbf{Diga:} ``Una derivada parcial responde una pregunta muy concreta:
¿qué ocurre con el costo si cambio una sola decisión y mantengo la otra fija?''

\textbf{Escriba primero el experimento de publicidad:}
\[
\frac{\partial J}{\partial x}
\quad\Longleftrightarrow\quad
\text{cambiar un poco }x\text{ mientras dejamos }y\text{ quieta}.
\]
Para nuestro modelo, al mirar solo los términos que contienen $x$:
\[
\frac{\partial}{\partial x}(x-12)^2=2(x-12),
\qquad
\frac{\partial}{\partial x}\bigl(2(y-8)^2\bigr)=0.
\]
El segundo término vale cero porque, en este experimento, $y$ no cambia.
Por eso:
\[
\boxed{\frac{\partial J}{\partial x}=2(x-12)}.
\]

\textbf{Repita el experimento con inventario:}
\[
\frac{\partial J}{\partial y}
\quad\Longleftrightarrow\quad
\text{cambiar un poco }y\text{ mientras dejamos }x\text{ quieta},
\]
\[
\frac{\partial}{\partial y}\bigl(2(y-8)^2\bigr)=4(y-8),
\qquad
\boxed{\frac{\partial J}{\partial y}=4(y-8)}.
\]

\textbf{Ejemplo hablado:} en $(x,y)=(2,2)$,
\[
\frac{\partial J}{\partial x}=-20,\qquad
\frac{\partial J}{\partial y}=-24.
\]
El primer número dice ``aumentar $x$ reduce el costo localmente'' y el segundo
 dice lo mismo para $y$. Juntos forman el gradiente:
\[
\boxed{\nabla J(2,2)=(-20,-24)}.
\]

\textbf{Compruebe el signo:} señal negativa $\Rightarrow$ al restarla aumentamos esa
variable; señal positiva $\Rightarrow$ al restarla la reducimos. No se requiere
demostrar reglas de derivación: basta entender qué decisión se mantiene fija y qué
decisión se está probando.

\section*{Ruta general del algoritmo}
\textbf{Escriba y deje visible la fórmula general:}
\[
\boxed{\boldsymbol{\theta}_{k+1}
=\boldsymbol{\theta}_k-\eta\,\nabla J(\boldsymbol{\theta}_k)}
\]
En dos variables, $\boldsymbol{\theta}_k=(x_k,y_k)$:
\[
\boxed{
\begin{aligned}
x_{k+1}&=x_k-\eta\,\frac{\partial J}{\partial x}(x_k,y_k),\\
y_{k+1}&=y_k-\eta\,\frac{\partial J}{\partial y}(x_k,y_k).
\end{aligned}}
\]

\textbf{Escriba la ruta de aplicación:}
\begin{enumerate}[leftmargin=1.8em]
\item Definir las decisiones $x$ e $y$.
\item Definir la función de costo, error o pérdida $J(x,y)$.
\item Elegir $(x_0,y_0)$ y el tamaño del paso $\eta$.
\item Calcular las dos derivadas parciales.
\item Actualizar simultáneamente $x$ e $y$.
\item Comprobar si el costo disminuyó.
\item Repetir hasta que la mejora sea menor que una tolerancia
      $\varepsilon$ o se alcance el máximo de rondas $K_{\max}$.
\end{enumerate}

\[
\boxed{\text{ruta}=(x_0,y_0)\longrightarrow(x_1,y_1)
\longrightarrow(x_2,y_2)\longrightarrow\cdots}
\]
\textbf{Diga:} ``Cada punto es una política completa y cada flecha es una
corrección propuesta por el gradiente.''

\section*{La regla de parada --- ¿cuándo termina?}
\textbf{Escriba antes de iniciar las rondas:}
\[
\boxed{\varepsilon=0.01\qquad K_{\max}=30}
\]
\[
\boxed{
\text{DETENER si }
J(x_k,y_k)-J(x_{k+1},y_{k+1})<\varepsilon
\quad\text{o}\quad k+1=K_{\max}}
\]

\textbf{Diga:} ``Después de cada ronda medimos cuánto bajó el costo. Terminamos
si la mejora ya es demasiado pequeña o si agotamos el número máximo de rondas.
El algoritmo no necesita llegar exactamente al punto ideal.''

\textbf{Final del ejemplo con $\eta=0.10$:}
\[
\begin{array}{c|c|c|c}
\text{ronda}&(x_k,y_k)&J(x_k,y_k)&\text{mejora}\\
\hline
19&(11.856,\ 8.000)&100.0208&0.0117\\
\mathbf{20}&\mathbf{(11.885,\ 8.000)}&\mathbf{100.0133}&
\mathbf{0.0075<0.01}
\end{array}
\]
\[
\boxed{\text{FIN en la ronda 20: la mejora es menor que }\varepsilon=0.01}
\]

\textbf{Aclaración para la demostración corta:} si solo se dibujan ocho rondas,
escriba:
\[
\boxed{\text{FIN de la demostración en la ronda 8: se alcanzó el límite elegido}}
\]
\[
\boxed{\text{El algoritmo podría continuar}}
\]

\section*{Tablero 4 --- Datos de la empresa (40--43 min)}
\[
\boxed{(x_0,y_0)=(2,2)}\qquad\boxed{\eta=0.10}
\]
\textbf{Diga:} ``La empresa empieza con poca publicidad y poco inventario. Veamos
qué recomienda el mapa.''

\newpage
\section*{Tablero 5 --- Primera ronda, completa y lenta (43--55 min)}
\textbf{Escriba una línea, pregunte el signo y recién después escriba la siguiente:}
\[
\begin{aligned}
g_x&=2(2-12)=\boxed{-20},&
g_y&=4(2-8)=\boxed{-24}.\\
x_1&=2-0.10(-20)=\boxed{4},&
y_1&=2-0.10(-24)=\boxed{4.4}.
\end{aligned}
\]

\textbf{Interpretación:} ambas señales negativas; aumentamos publicidad e inventario.

\textbf{Costo antes:}
\[
J(2,2)=(-10)^2+2(-6)^2+100=100+72+100=\boxed{272}.
\]
\textbf{Costo después:}
\[
J(4,4.4)=(-8)^2+2(-3.6)^2+100
=64+25.92+100=\boxed{189.92}.
\]
\[
\boxed{\text{reducción}=272-189.92=82.08}
\]

\textbf{Aclare:} ``Las dos cuentas de actualización ocurren en la misma ronda. No
usamos $g_x$ para cambiar $y$ ni $g_y$ para cambiar $x$.''

\newpage
\section*{Tablero 6 --- Segunda y tercera ronda (55--68 min)}
\textbf{Segunda ronda:}
\[
\begin{aligned}
g_x&=2(4-12)=-16,&g_y&=4(4.4-8)=-14.4,\\
x_2&=4-0.10(-16)=\boxed{5.6},&
y_2&=4.4-0.10(-14.4)=\boxed{5.84}.
\end{aligned}
\]
\[
J(5.6,5.84)=(-6.4)^2+2(-2.16)^2+100
=40.96+9.3312+100=\boxed{150.2912}.
\]

\textbf{Tercera ronda: pida que una pareja la complete:}
\[
\begin{aligned}
g_x&=2(5.6-12)=\underline{\hspace{1.8cm}},&
g_y&=4(5.84-8)=\underline{\hspace{1.8cm}},\\
x_3&=5.6-0.10(\underline{\hspace{1.3cm}})=\underline{\hspace{1.5cm}},&
y_3&=5.84-0.10(\underline{\hspace{1.3cm}})=\underline{\hspace{1.5cm}}.
\end{aligned}
\]
\textbf{Solución para revelar:}
\[
g_x=-12.8,\quad g_y=-8.64,\quad
(x_3,y_3)=(\boxed{6.88},\boxed{6.704}).
\]
\textbf{Costo de la tercera ronda:}
\[
J(6.88,6.704)=(-5.12)^2+2(-1.296)^2+100
=26.2144+3.359232+100=\boxed{129.573632}.
\]

\textbf{Pregunte:} ``¿Por qué las correcciones se hacen menores?'' Respuesta:
las señales se acercan a cero al acercarnos a $(12,8)$.

\section*{Tabla de seguimiento para dejar visible}
\small
\begin{center}
\begin{tabular}{r|rr|rr|r}
\toprule
$k$&$x_k$&$y_k$&$g_x$&$g_y$&$J(x_k,y_k)$\\
\midrule
0&2.000&2.000&-20.000&-24.000&272.000\\
1&4.000&4.400&-16.000&-14.400&189.920\\
2&5.600&5.840&-12.800&-8.640&150.291\\
3&6.880&6.704&-10.240&-5.184&129.574\\
4&7.904&7.222&-8.192&-3.110&117.987\\
\bottomrule
\end{tabular}
\end{center}
\normalsize

\newpage
\section*{Animaciones del cuaderno (68--90 min)}
\textbf{Antes de abrir las celdas de animación, escriba:}
\[
\text{predicción: después de tres rondas }(x,y)\approx
(\underline{\hspace{2cm}},\underline{\hspace{2cm}})
\]

\textbf{Animación 1, mapa:} ``El punto rojo atraviesa curvas de nivel cada vez más
interiores. Eso significa que el costo baja.''

\textbf{Pregunte:}
\begin{enumerate}
\item ¿Qué coordenada cambia más al principio?
\item ¿Por qué el camino no es recto?
\item ¿Qué representa el centro verde?
\end{enumerate}

\textbf{Animación 2, superficie 3D:} ``Ahora vemos la misma trayectoria como una
montaña. La altura es el costo. El punto no solo se acerca al centro: también baja.''

\textbf{Si no se reproduce en pantalla:} ejecutar de nuevo la celda, abrir el
cuaderno en Jupyter/VS Code y comprobar que aparece el GIF incrustado. La salida
incluye texto con la ronda, $(x,y)$ y el costo para seguir la secuencia sin animación.

\section*{Tablero 7 --- Tamaño del paso (90--105 min)}
\[
\begin{array}{c|c|c}
\eta=0.03&\eta=0.10&\eta=0.30\\
\text{lento}&\text{estable}&\text{agresivo}\\
\text{muchos intentos}&\text{avance prudente}&\text{puede oscilar}
\end{array}
\]

\textbf{Diga:} ``Una empresa puede sobrecorregir: demasiado inventario hoy y muy
poco mañana. El tamaño del paso representa prudencia.''

\textbf{Regla de diagnóstico:}
\[
\begin{array}{ll}
\text{costo baja establemente}&\Rightarrow\text{continuar}\\
\text{costo baja casi nada}&\Rightarrow\text{paso quizá muy pequeño}\\
\text{costo alterna o crece}&\Rightarrow\text{reducir el paso}
\end{array}
\]

\newpage
\section*{Tablero 8 --- Actividad: comité empresarial (105--115 min)}
\textbf{Escriba:}
\[
(x,y)=(18,3),\qquad \eta=0.10.
\]
\[
\begin{aligned}
g_x&=2(18-12)=\underline{\hspace{1.5cm}},&
g_y&=4(3-8)=\underline{\hspace{1.5cm}},\\
x_{\rm nuevo}&=18-0.10(g_x)=\underline{\hspace{1.5cm}},&
y_{\rm nuevo}&=3-0.10(g_y)=\underline{\hspace{1.5cm}}.
\end{aligned}
\]

\textbf{Solución:}
\[
g_x=12,\quad g_y=-20,\quad
(x_{\rm nuevo},y_{\rm nuevo})=(16.8,5).
\]
\[
J(18,3)=36+50+100=\boxed{186}
\]
\[
J(16.8,5)=23.04+18+100=\boxed{141.04};
\qquad\boxed{\text{reducción}=44.96}.
\]

\textbf{Recomendación modelo:} ``Reducir publicidad de 18 a 16.8 y aumentar
inventario de 3 a 5; el costo estimado baja 44.96 unidades.''

\section*{Tablero 9 --- Cierre (115--120 min)}
\[
\boxed{\text{dos decisiones}\quad\text{dos señales}\quad\text{un paso conjunto}}
\]
\begin{enumerate}
\item El mapa muestra todas las políticas posibles.
\item $g_x$ indica cómo corregir publicidad.
\item $g_y$ indica cómo corregir inventario.
\item Cada variable se actualiza en su propia ecuación.
\item Repetimos y verificamos que el costo baje.
\item Terminamos cuando la mejora es menor que $\varepsilon$ o alcanzamos
      $K_{\max}$.
\end{enumerate}

\textbf{Boleto de salida:}
\[
\boxed{\text{En este caso }g_x\text{ significa }\underline{\hspace{3cm}}
\text{ y }g_y\text{ significa }\underline{\hspace{3cm}}.}
\]

\section*{Lista de control}
\begin{itemize}[leftmargin=1.5em]
\item Ejecutar las celdas en orden para que las dos animaciones aparezcan.
\item Verificar que la salida muestra un GIF, no solo el código.
\item Explicar derivadas parciales como ``cambiar una decisión y mantener la otra fija''.
\item No exigir una demostración formal: concentrarse en dirección, cuentas e interpretación.
\item Mantener la distinción entre costo total $J$ y decisiones $(x,y)$.
\item Escribir la regla de parada antes de comenzar y anunciar explícitamente
      por qué terminó el proceso.
\end{itemize}
\end{document}
